% Folder 156-8, batch 06 — archivists' pages 101–120 (author's pages 93–112).
% Pass: Opus 5.5 (claude-opus-5-5), 2026-10-02 — first pass, unchecked against the pages by a human.
%
% Notation fixed for this batch: his underlined script capital (a Phi-like
% letter) is set $\Phi$; $\mathcal{M}$, $\mathcal{M}_0$, $\mathcal{M}_1$,
% $\mathcal{L}$; the relation he writes |o| is set $|{\circ}|$; the order
% "< pos" is set $\underset{\mathrm{pos}}{<}$; "<" with a 1 above is set
% $\overset{1}{<}$; the incidence triangle is set $\triangleleft$.
\documentclass[11pt,a4paper]{article}
\input{../preamble/grothendieck}

\folder{156-8}
\batch{6}
\pages{101}{120}
\foldertitle{[Chapitre] VIII. Analysis situs (quatrième mouture) : notes manuscrites (26/06-04/07/1986).}
\dating{1986}
\watermark{Édition de démonstration}

\begin{document}

\page{101}\note{p. 93 de l'auteur. La page s'ouvre au milieu d'une démonstration commencée dans le lot précédent : il s'agit de montrer que $\operatorname{supp}^\circ(P_a)=\{P_a\}$, puis (Cor.~2) l'énoncé analogue pour $S_{a,b}$.}
distinct de $P_a$ est $|{\circ}|$ à $P_a$), d'où résulte \uncertain{aussitôt} (à cause de l'\uncertain{antiréflexivité} de $|{\circ}|$) que $\operatorname{supp}^\circ(P_a)=\{P_a\}$. Si $a$ n'est ni plus petit ni plus grand élément, soient \struck{$x<$}
\[ x<a<y . \]
Alors \struck{$\xi$} $S_{x,a}\in\operatorname{cosupp}^\circ(P_a)$, donc $X\,|{\circ}|\,S_{x,a}$, donc $\delta X\geq a$ ou $\delta X\leq x$. Mais on ne peut avoir $\delta X\leq x$, car on aurait $X\,|{\circ}|\,a$, i.e. $X\in\operatorname{cosupp}^\circ(P_a)$, ce qui contredit $X\in\operatorname{supp}^\circ(P_a)$ (\uncertain{ou l'antiréflexivité}) : $\operatorname{cosupp}^\circ(\Phi)\cap\operatorname{supp}^\circ(\Phi)=\emptyset$, pour toute partie $\Phi$ de $\mathcal{M}$). Donc $\delta X\geq a$. L'\struck{\ill{}} argument symétrique, utilisant $S_{a,y}$ au lieu de $S_{x,a}$, donne $\delta X\leq a$, d'où la conclusion.

\textbf{Cor. 2} Raisonnement analogue ; \struck{on} \struck{si $a$ est un p\ill{}} prouver successivement
\[ \delta X\geq a, \quad\text{puis}\quad \delta X\leq b \]
en distinguant, pour la première, le cas où $a$ est un plus petit élément (\uncertain{facile}) et le cas où il ne l'est pas, donc $\exists x$,
\[ x<a<b \]
et on aurait $S_{x,a}\in\operatorname{cosupp}^\circ(S_{a,b})$, donc $X\,|{\circ}|\,S_{x,a}$, donc $\delta X\leq x$ ou $\delta X\geq a$, et la première \ill{} \uncertain{exclue} comme \ill{} devant \ill{} avec $X\,|{\circ}|\,\widetilde{S}_{a,b}$).

\page{102}\note{p. 94 de l'auteur.}
\uncertain{Ceci} \ill{} \ill{} fait par l'hyp. \uncertain{par exemple}

\begin{tikzcd}[column sep=small, row sep=small]
 & b & \\
a \arrow[ur] & & c \arrow[ul]
\end{tikzcd}

\ill{} \[ = \operatorname{cosupp}^\circ(P_b) = \mathcal{M}\setminus\{P_b\} \]
\[ \operatorname{cosupp}^\circ(\{P_a\}) = \{P_b\text{\struck{$\}$}},\ S_{a,b},\ S_{c,b} \]
\[ \operatorname{supp}^\circ(\{P_a\}) = \{P_a,\ P_b\} \neq \{P_a\} \]
\note{les termes $S_{a,b}, S_{c,b}$ de la première ligne sont reliés par un trait à l'accolade de la seconde.}
\[ \operatorname{cosupp}^\circ(\widetilde{S}_{a,b}) = \emptyset \]
\[ \operatorname{supp}^\circ(\widetilde{S}_{a,b}) = \mathcal{M} \neq \operatorname{Omb}(S_{a,b}) \]
(NB $\mathcal{M}=\operatorname{Omb}(S_{a,b})\amalg\{S_{c,b},P_c\}$, ici $\operatorname{card}\mathcal{M}=5$)

\marginal{\uncertain{Canular}}

Pour \uncertain{contre-exemple} \uncertain{pour} Pb \ill{}

\begin{tikzcd}[column sep=small, row sep=small]
 & b & \\
a \arrow[ur] & & d \arrow[ul] \\
 & c \arrow[ul] \arrow[ur] &
\end{tikzcd}

p.ex.

\struck{$\operatorname{cosupp}^\circ P_a = \{P_b, P_c\}, S_{a,b}, S_{c,a}\}$}

\struck{$\operatorname{supp}^\circ P_a =$}

\note{le lemme qui suit est biffé de trois traits obliques.}
\struck{\textbf{Lemme 2}} \struck{Soient} \struck{$\Phi\subset\mathcal{M}$, $\Psi\subset\mathcal{M}$} \struck{$X,Y\in\mathcal{M}$} \struck{\ill{} $X\,\|\,Y$, et soit $\Phi\subset\operatorname{Omb}(X)$ ($=\operatorname{supp}\widetilde{X}$ par lemme 1), $\Psi\subset\operatorname{Omb}(Y)$ ($=\operatorname{supp}\widetilde{Y}$).} \struck{Alors}
\[ \text{\struck{$\operatorname{supp}^\circ(\Phi\cup\Psi)=\operatorname{supp}^\circ\Phi\cup\operatorname{supp}^\circ\Psi$.}} \]

\marginal{4 juillet}
Détermination des profigures, ou ensembles constructibles dans $\mathcal{M}$.

\struck{Les} P\ill{} la relation d'ordre réduite, une profigure se décompose en composantes connexes \add{(profigures connexes)}. Les composantes connexes \uncertain{sont} des types suivants
\struck{a) \ill{}} \struck{$\Phi=\{P_a\}$} \struck{(c'est une} figure fermée)
\struck{b) \ill{}} \struck{b) Figure}

\page{103}\note{p. 95 de l'auteur.}
a) \struck{Pré}figures fermées connexes \add{$=$ figures connexes} ($=$ déterminées par les éléments de $\operatorname{Drap}^*(\mathcal{L})$, cf. plus haut)

b) Figures \uncertain{fermantes} connexes : \struck{cas où} elles sont données par $T\in\operatorname{Drap}^*(\mathcal{L})$, \uncertain{$\operatorname{card}T\geq 2$}, et une partie de
\[ \partial F_T \overset{\text{déf}}{=} \partial T = \{\text{ensemble formé du plus petit et du plus grand élément de } T=F_0\} \]
\[ \Phi = F\setminus\partial F \]
NB on a fini par identifier les $P_a$ aux $a$, i.e. $\mathcal{M}_0$ à $\mathcal{L}$…

c) Cas général : on se donne une partie $\alpha$ de $\partial F$, et on prend
\[ \Phi = F\setminus\alpha \]

\textbf{Ex} Segment fermé, semi-ouvert, ouvert (cas $\operatorname{card}T=2$)

Soit $\text{Préfigcomm}(\mathcal{M})$ l'ensemble des préfigures connexes, \uncertain{qui sont} \ill{} d'\uncertain{explicitation}, identifiées aux couples $(T,\alpha)$, $T\in\operatorname{Drap}^*(\mathcal{L})$, $\alpha\in\partial T$ (\struck{frontière} « bord » de la figure \add{connexe} $F_T$ ; NB si $F_T=\{P_\alpha\}$, on \add{prendra} $\partial F_T=\emptyset$).

\marginal{On pose $F=\overline{\Phi}$ (plus petite figure \struck{\ill{}} \uncertain{contenant} $\Phi$), et $\partial\Phi=\partial F$}

On va définir sur
\[ \text{Préfigcomm}(\mathcal{M}) \]
une relation d'ordre \struck{$<$} en écrivant
\[ \Phi < \Phi' \iff
\begin{cases}
\text{pour tout } X\in\Phi,\ X'\in\Phi',\ \text{on a } X \underset{\mathrm{pos}}{<} X' \\
\text{a) [\dots]}\ \operatorname{card}(\delta\Phi\cap\delta\Phi')\leq 1 \\
\text{b) si } x\in\partial\Phi\cap\partial\Phi',\ \text{on a } x\notin\Phi,\ x\notin\Phi'
\end{cases} \]
\note{après $X \underset{\mathrm{pos}}{<} X'$, les mots « au sens \ill{} » sont biffés, puis « \uncertain{plus bas} ».}

NB b) indique que si $\Phi$ \uncertain{de dim} $0$, i.e. $\Phi=\{P_a\}$, alors \ill{}

\page{104}\note{p. 96 de l'auteur.}
Soit $\Phi$ une profigure, $F$ la figure engendrée (on pose $F=\overline{\Phi}$, c'est bien l'adhérence de $\Phi$ pour $\leq$ sur $\mathcal{M}$), on pose
\[ \delta\Phi \overset{\text{déf}}{=} \delta F = F\cap\mathcal{M}_0
= \{a\in\mathcal{L} \mid P_a\in F,\ \text{i.e. } \exists X\in\Phi,\ P_a\leq X\} \]
\[ \text{\struck{$\partial\Phi \overset{\text{déf}}{=} \partial F =$}}\quad
\text{\struck{$\emptyset$ si $\operatorname{card}\delta F\leq 1$ ; ensemble formé des plus petit et plus grand}} \]
\note{au-dessus de la ligne biffée : « “extrémités” de $\delta F$ ».}

NB $\delta\Phi=\delta F\in\operatorname{Drap}^*(\mathcal{L})$, \uncertain{ens.} des \struck{\ill{}} \struck{pour $\Phi$, i.e. pour $F$} de la profigure $\Phi$. Il est vide ssi $\Phi=\emptyset$, de card $1$ ssi $\Phi=$ \struck{\ill{}} $\{P_a\}$, auquel cas $\Phi=F$.

\marginal{Montrer que si $a$ \ill{} de $\Phi$, \ill{} \ill{} \ill{} \ill{} $P_a\in\Phi$, ex : $\Phi=\{S_{a,b}\}$}

\textbf{Cas $\Phi$ connexe}
\[ \partial\Phi \overset{\text{déf}}{=}
\begin{cases}
\text{a) } \emptyset \text{ si } \operatorname{card}\delta\Phi\leq 1, \text{ i.e. } \Phi=\emptyset \text{ ou } \{P_a\} \\
\text{b) bord de } \delta\Phi \text{ dès } \operatorname{card}\delta\Phi\geq 2
\end{cases} \]
i.e. ens. formé des extrémités du drapeau $\delta\Phi$ (plus petit et plus grand élément)

\marginal{bord de $\Phi$ : $\partial\Phi=\partial F$}

ainsi, $\partial\Phi=\emptyset$ ou $\in\operatorname{Drap}_2(\mathcal{L})$.

[Cas $\Phi$ connexe
\[ \Phi^{\bullet} =
\begin{cases}
\partial\Phi & \text{si } \operatorname{card}\delta\Phi\geq 2 \\
\{a\} & \text{si } \Phi=\{P_a\}, \text{ i.e. } \operatorname{card}\delta\Phi=1 \\
\emptyset & \text{si } \Phi=\emptyset, \text{ i.e. } \operatorname{card}\delta\Phi=0
\end{cases} \]
]

$\Phi$ pas néc. connexe
\[ \partial\Phi = \bigcup \partial\Phi_i \]
\[ \Phi^{\bullet} = \bigcup \Phi_i^{\bullet} \]
où les $\Phi_i$ sont les comp. connexes.

\textbf{Ex 1} \note{le chiffre est douteux.}
\[ \Phi = \{S_{a,b},\ S_{b,c},\ P_d\} \]
où $a<b<c<d$, comp. connexes $S_{a,b}$, $S_{b,c}$, $P_d$.
\[ \delta\Phi = \{a,b,c,d\} \]
\[ \partial\Phi = \{a,b,c\} \]
\[ \Phi^{\bullet} = \{a,b,c,d\} \]

\marginal{\uncertain{définition} \ill{} plus petit \uncertain{ou} plus grand \ill{} — mais on n'en a pas \uncertain{besoin}… On pose, pour $\Psi\in\mathcal{M}$, $a\in\mathcal{L}$, \add{profigure}
\[ \operatorname{ordre}(a,\Psi) = \operatorname{card}(\{X\in\mathcal{M}_1\cap\Psi \mid a\triangleleft X\}) \]
(cet ordre est $0$, $1$ ou $2$) : $0$ point extérieur (si $a\notin\Psi$) ou isolé (si $a\in\Psi$) ; $1$ point bord ; $2$ point frontière intérieur (si $a\notin\Psi$) ou point \struck{\ill{}} redondant (\uncertain{ou} de morcellement) si $a\in\Psi$
\[ \partial\Phi = \{a\in\mathcal{L} \mid \operatorname{ordre}(a,\overline{\Phi})=1\} \]}
\note{la note marginale, écrite en biais, porte plusieurs biffures dans la formule de $\partial\Phi$.}

\page{105}\note{p. 97 de l'auteur.}
\textbf{Ex 3}
\[ \Phi = \{S_{a,b},\ P_b,\ S_{b,c},\ P_d\}, \]
comp. connexes $\{S_{a,b}, P_b, S_{b,c}\}$ et $P_d$
\[ \delta\Phi = \{a,b,c,d\} \]
\[ \partial\Phi = \{a,c\} \]
\[ \Phi^{\bullet} = \{a,c,d\} \]
\uncertain{Tous les termes diffèrent.}

NB Dans tous les cas, on a
\[ \partial\Phi \subset \Phi^{\bullet} \subset \delta\Phi \]
\note{sous la première inclusion, un mot souligné, \uncertain{connexité}, renvoyant à : « (si $\exists$ plus grand ou plus petit élément dans $\mathcal{L}$ !) ».}

Et \struck{$\Phi^{\bullet}$ \ill{}} On \ill{}
\[ \partial\Phi = \{a\in\delta\Phi \mid \exists X\in\Phi \text{ tel que } P_a<X, \text{ et de plus, quand } P_a\in\Phi, \text{ cet } X \text{ est unique}\} \]
\note{entre les lignes de l'accolade, des essais biffés : « ou bien $P_a\notin\Phi$, mais existe $X\in\Phi$ ou… », « $P_a<X$ », « ou bien $P_a\in\Phi$, \ill{} ».}
\[ = \text{ens. des sommets non isolés de } \Phi, \]
qui, de plus, quand ils appartiennent à $\Phi$, ne sont pas « intérieurs » à $\Phi$.

\marginal{NB Dire que si $P_a\in\Phi$, il existe au plus deux $X\in\Phi$ tels que $P_a<X$}

\[ \Phi^{\bullet} = \partial\Phi \amalg \Phi_{\mathrm{isol}} \]
où $\Phi_{\mathrm{isol}}$ est l'ens. des \struck{points} sommets isolés de $\Phi$.
\uncertain{canular}…

Dans $\mathcal{M}$, on va introduire une relation d'ordre, en écrivant
\[ X \underset{\mathrm{pos}}{<} Y \iff
\begin{cases}
\text{a) } X=P_x,\ Y=P_y & x<y \\
\text{b) } X=P_x,\ Y=S_{a,b} & x\leq a<b \\
\text{c) } X=S_{a,b},\ Y=P_y & a<b\leq y \\
\text{d) } X=S_{a,b},\ Y=S_{a',b'} & a<b\leq a'<b'
\end{cases} \]
\[ X \underset{\mathrm{pos}}{<} Y \iff X\,|{\circ}|\,Y \text{ et } \operatorname{or}(\delta X) < \operatorname{ex}(\delta Y) \]
\note{sous $\operatorname{or}(\delta X)$ : « \struck{plus petit} élément de $\delta X$ » ; sous $\operatorname{ex}(\delta Y)$ : « plus grand él. de $\delta Y$ ».}

\marginal{NB l'indice pos dans $\underset{\mathrm{pos}}{<}$ \uncertain{est} \ill{}, \ill{} \uncertain{pas besoin de} \ill{} $<,\leq$ pour la relation \uncertain{d'incidence} dans $\mathcal{M}$…}

\page{106}\note{p. 98 de l'auteur.}
En d'autres termes, on a trouvé une relation d'ordre sur $\mathcal{M}$, telle que la relation $X\,|{\circ}|\,Y$ sur $\mathcal{M}$ soit la relation de disjonction \struck{\ill{}} associée
\[ X\,|{\circ}|\,Y \iff X \underset{\mathrm{pos}}{<} Y \ \text{ou}\ Y \underset{\mathrm{pos}}{<} X \]
Si $\Phi,\Psi$ sont deux parties quelconques de $\mathcal{M}$, on pose
\[ \Phi \underset{\mathrm{pos}}{<} \Psi \overset{\text{déf}}{\iff} \forall X\in\Phi,\ Y\in\Psi,\quad X<Y . \]

NB Si $\Phi,\Psi$ \uncertain{sont} des profigures (\uncertain{et non} $\Phi \underset{\mathrm{pos}}{\leq} \Psi$), \uncertain{la} $\Phi\cup\Psi$ \ill{} relation d'ordre sur \ill{} précédemment, \ill{} des \uncertain{drapeaux} comme \uncertain{ensembles} des parties de $\mathcal{M}_0\subset\mathcal{M}$. \struck{Ceci}

Dans le cas où $\Phi,\Psi$, \uncertain{et} des \uncertain{profigures} admissibles, on va introduire une relation un peu plus forte
\[ \Phi \underset{\mathrm{pos}}{\overset{1}{<}} \Psi \overset{\text{déf}}{\iff}
\begin{cases}
\Phi \underset{\mathrm{pos}}{<} \Psi, \text{ et dans le cas où} \\
\operatorname{ex}\Phi = \operatorname{or}\Psi \ (\text{\uncertain{noté} } a), \text{ on a} \\
P_a\notin\Phi,\ P_a\notin\Psi
\end{cases} \]
\note{au-dessus de « dans le cas où », une addition peu lisible, \uncertain{« dans ce cas »}.}

i.e. on exclut un cas comme
\drawing{deux segments consécutifs, $\Phi$ à gauche et $\Psi$ à droite, raccordés en un point marqué $(\bullet)$ : « …——(•)——… » ou l'inverse.}

\marginal{NB si $\Phi \underset{\mathrm{pos}}{<} \Psi$ \add{profigures}, alors \struck{$\delta\Phi \underset{\mathrm{pos}}{<} \delta\Psi$ ou $\delta\Phi\cap\delta\Psi=\{a\}$, $\operatorname{ex}(\delta\Phi)=\operatorname{or}(\delta\Psi)$} \struck{alors $P_a\notin\Phi\cap\Psi$} $\delta\Phi\cap\delta\Psi=\emptyset$ ou de \uncertain{cardinal} $1$, $a\in\delta\Phi\cap\delta\Psi$, et si $\Phi\cap\Psi=\emptyset$, et si $a=\operatorname{ex}(\delta\Phi)=\operatorname{or}(\delta\Psi)$, on a :
\[ \Phi \underset{\mathrm{pos}}{\leq} \Psi \iff \operatorname{ex}\Phi < \operatorname{or}\Psi, \]
ou $\operatorname{ex}\Phi=\operatorname{or}\Psi$ et $P_a\notin\Phi\cap\Psi$}
\note{note marginale en biais, en partie biffée ; la lecture est incertaine dans le détail.}

\page{107}\note{p. 99 de l'auteur.}
Cas du « raccord », \struck{quand $\Phi$} En d'autres termes, on exclut le cas où $\Phi,\Psi$ seraient \add{\uncertain{ils pas}} des parties \uncertain{ouvertes} et fermées de $\Phi\cup\Psi$.

Ceci posé, voici la structure d'une profigure quelconque, en termes de profigures connexes.

\struck{\ill{}} Soit
\[ D = \{\Phi_1 \underset{\mathrm{pos}}{\overset{1}{<}} \Phi_2 \underset{\mathrm{pos}}{\overset{1}{<}} \cdots \underset{\mathrm{pos}}{\overset{1}{<}} \Phi_n\} \]
un drapeau \struck{pour} dans $\text{Préfigcomm}(\mathcal{M})$
\[ D \in \operatorname{Drap}^*_{\overset{1}{<}}(\text{Préfigcomm}(\mathcal{M})) \]
on lui associe la profigure
\[ \Phi_D = \bigcup_{\Phi\in D}\Phi \ \text{\struck{$\ill{}$}}\ = \bigcup_{i\in[1,n]}\Phi_i \]
(réunion disjointe). \uncertain{Ce} fait qu'il fallait pour que les $\Phi_i$ soient les comp. connexes de $\Phi$. On trouve \ill{} que dans le cas général \struck{\ill{}}, \struck{\ill{}} les figures \ill{} \uncertain{intérieures} \ill{}, si

\textbf{Lemme} Soient $\Phi,\Psi$ profigures connexes. Pour que $\Phi\cap\Psi=\emptyset$ et $\Phi\cup\Psi$ soit une profigure, \struck{\ill{} $\Phi \underset{\mathrm{pre}}{\leq} \Psi$} il faut et il suffit que
\[ \Phi \underset{\mathrm{pos}}{\leq} \Psi . \]
\struck{Pour que $\Phi$} \struck{\ill{}}

\page{108}\note{p. 100 de l'auteur.}
et $\Psi$ soient les comp. connexes de $\Phi\cup\Psi$, i.e. $\Phi\cup\Psi$ déconnecté, il f. et s.
\[ \text{que}\quad \Phi \overset{1}{<} \Psi \ \text{ou}\ \Psi \overset{1}{<} \Phi, \ \text{i.e.} \]
\[ \{\Phi,\Psi\} \in \operatorname{Drap}_{\overset{1}{<}}(\text{Préfigcomm}(\mathcal{M})) . \]

En résumé, on trouve donc la petite proposition récapitulative

\textbf{Théorème} 1) Pour tout $D\in\operatorname{Drap}^*(\mathcal{L})$, et toute partie $\alpha$ de \struck{$\delta$}$\partial D$ (cf. …), soit
\[ \Phi_{D,\alpha} = F_D\setminus\alpha , \]
où $F_D$ est la figure fermée des \uncertain{strates} \uncertain{constructibles} du drapeau $D$. On trouve ainsi une bijection
\[ (D,\alpha) \longmapsto \Phi_{D,\alpha} \]
de l'ens. des paires $D,\alpha$ sur l'ens. des profigures connexes.

\marginal{$\partial D = \begin{cases} \emptyset & \text{si } \operatorname{card}D=1 \\ \{\operatorname{or}D,\ \operatorname{ex}D\} & \text{si } \operatorname{card}D\geq 2 \end{cases}$}

2) Pour que $\Phi$ soit une figure, i.e. fermée, il f. et s. que $\alpha=\emptyset$. Pour que $\Phi$ soit ouverte, \struck{\ill{}} il f. et s., \uncertain{sous} \uncertain{réserve de} Hyp \struck{\ill{}} et en supposant que $\mathcal{L}$ \uncertain{n'a} pas de plus petit ni de plus grand élément, que $\operatorname{card}D\geq 2$ et $\alpha=\partial D$. En \struck{dirigeant \ill{}} \uncertain{éliminant} \uncertain{cette} Hyp. et \uncertain{dans} \uncertain{le cas} où il existe un plus petit ou un

\page{109}\note{p. 101 de l'auteur.}
plus grand élément dans $\mathcal{L}$, soit \struck{$\xi$} $\Theta$ l'ens. \add{de la partie} de $\mathcal{L}$, de cardinal $\{0,1,2\}$, formée de ces éléments. Alors $\Phi$ est ouvert ssi
\[ \text{\struck{$\alpha=\partial D$}}\quad \operatorname{card}D\geq 2 \ \text{et}\ \alpha = \partial D - (\partial D\cap\Theta) . \]

3) Pour \struck{toute suite de profigures} \add{drapeau} \struck{connexes},
\[ D = \{\Phi_1 \overset{1}{<} \Phi_2 \overset{1}{<} \cdots \overset{1}{<} \Phi_n\} \]
dans l'ens. ordonné des \add{pro}figures connexes,
\[ D \in \operatorname{Drap}^*_{\overset{1}{<}}(\text{Préfigcomm}(\mathcal{M})), \]
soit
\[ \Phi_D = \bigcup_{\Phi\in D}\Phi = \bigcup_{1\leq i\leq n}\Phi_i . \]
C'est une réunion disjointe, et \struck{\ill{}} $\Phi_D$ est une profigure, admettant les $\Phi\in D$ comme composantes connexes.
L'application
\[ D \longmapsto \Phi_D = \bigcup_{\Phi\in D}\Phi \]
\[ \operatorname{Drap}^*_{\overset{1}{<}}(\text{Préfigcomm}(\mathcal{M})) \longrightarrow \text{Préfig}(\mathcal{M}) \]
est bijective. \struck{Pour}

4) Pour que $\Phi_D$ soit fermée, il f. et

\page{110}\note{p. 102 de l'auteur.}
suffit que les $\Phi\in D$ \uncertain{le soient} \uncertain{fermées}. Pour que $\Phi_D$ soit ouverte, il f. et s. que les $\Phi\in D$ le soient.

\marginal{Ces deux derniers \uncertain{sont} faits \ill{} \uncertain{triviaux}, \ill{} \ill{} \ill{} $\triangleq$ \ill{}}

Dans cette description des profigures et des figures, les relations $\overset{1}{<}$ n'ont pas intervenu encore, \uncertain{ni} (explicitement) la relation $|{\circ}|$ — à \uncertain{sa} place, c'est la relation plus explicite $\underset{\mathrm{pos}}{<}$ (\uncertain{ou plutôt} $<$) dans $\mathcal{M}$ qui nous a servi. Mais maintenant que je passe aux supports, il n'en sera pas de même.

Je \uncertain{veux} \struck{\ill{}} donner une détermination explicite du support d'une profigure quelconque. Cela se fait en plusieurs étapes

1°) Énoncé « réductif » : si \struck{les} $\Phi$ a pour composantes connexes les $\Phi_i$, on a
\[ \operatorname{supp}^\circ\Phi = \bigcup_i \operatorname{supp}^\circ(\Phi_i) \]
(réunion disjointe)

2°) Détermination de $\operatorname{supp}^\circ\Phi$ pour $\Phi$ profigure connexe. Pour ceci, trois étapes :

\page{111}\note{p. 103 de l'auteur.}
(a) Cas « irréductibles » des figures fermées élémentaires $\{P_a\}$, $\widetilde{S}_{a,b}=\{S_{a,b},P_a,P_b\}$, et des profigures associées $\widetilde{S}_{a,b}\setminus\alpha$ ($\alpha\in\{P_a,P_b\}$). Le résultat ici est connu (p.~91)

\marginal{\ill{} hyp. p.~91}

(i) Si $X\in\mathcal{M}_q$\note{l'indice $q$ est douteux.}, on a
\[ \operatorname{supp}X = \operatorname{supp}^\circ(\widetilde{X}) = \operatorname{Omb}(X) = \{Z\in\mathcal{M} \mid |Z|\subset|X|\} \]
\note{sous le premier membre : « déf. de $\operatorname{supp}X$ » ; sous le second signe $=$ : « th. » ; sous le troisième : « tautol. ».}
en particulier
\[ \operatorname{supp}\{P_a\} = \{P_a\} \]
\[ \operatorname{supp}\widetilde{S}_{a,b} = \left\{Z\in\mathcal{M} \;\middle|\;
\begin{array}{l} Z=P_x,\ a\leq x\leq b \\ Z=S_{a',b'},\ a\leq a'<b'\leq b \end{array}\right\} \]
(ii) Si \add{$X=S_{a,b}$ et} $\alpha\subset\{P_a,P_b\}$, i.e. $\alpha\subset\partial X$, on a
\[ \operatorname{supp}^\circ(\widetilde{X}\setminus\alpha) = \operatorname{supp}X\setminus\alpha \]
\[ = \left\{Z\in\mathcal{M} \;\middle|\;
\begin{array}{l} Z=P_x,\ a\leq x\leq b \text{ et } x\notin\alpha \\ Z=S_{a',b'},\ a\leq a'<b'\leq b \end{array}\right\} \]

Une autre façon \add{(plus élégante)} d'expliciter 2° est ainsi, \struck{$\forall X\in\mathcal{M}$}

$\alpha$) $\forall X\in\mathcal{M}$
\[ \operatorname{supp}^\circ X = \operatorname{Omb}^\circ X = \{Z\in\mathcal{M} \mid Z \overset{\circ}{<} X\} \]
en particulier
\[ \operatorname{supp}^\circ S_{a,b} = \left\{Z\in\mathcal{M} \;\middle|\;
\begin{array}{l} Z=P_x,\ a<x<b \\ Z=S_{a',b'},\ a\leq a'<b'\leq b \end{array}\right\} \]
$\beta$) Si $\beta\subset\partial S_{a,b}$, on a
\[ \operatorname{supp}^\circ(S_{a,b}\cup\beta) = \operatorname{supp}^\circ(S_{a,b})\cup\beta = \operatorname{Omb}^\circ(S_{a,b})\cup\beta \]
\[ = \left\{Z\in\mathcal{M} \;\middle|\;
\begin{array}{l} Z=P_x,\ a<x<b \text{ ou } x\in\beta \\ Z=S_{a',b'},\ a\leq a'<b'\leq b \end{array}\right\} . \]

\marginal{$\alpha$ et $\beta$ \uncertain{reliés} : $\beta=\partial X\setminus\alpha$, $\alpha=\partial X\setminus\beta$}

(b) Cas général : Considérons $T\in\operatorname{Drap}^*(\mathcal{L})$, $\operatorname{card}T\geq 3$, \struck{$\alpha$} $a=\operatorname{or}T$, $b=\operatorname{ex}(T)$, \struck{$\alpha$} $\alpha\subset\partial T=\{a,b\}$, $\beta$ son compl. dans $\partial T$, $\Phi=F_T\setminus\alpha$. On trouve alors :

\page{112}\note{p. 104 de l'auteur.}
\[ \operatorname{supp}\Phi = \operatorname{Omb}(\widetilde{S}_{a,b})\setminus\alpha = \operatorname{Omb}^\circ(S_{a,b})\cup\beta \]
\[ (\Phi = \Phi_{T,\alpha}) \qquad \text{(réunion disjointe)} \]

On a déjà fait (a) plus haut, il faut prouver encore (1°) et (b).

Prouvons 1°). Cela résulte, par réc. sur le nb de composantes connexes, du

\textbf{Lemme} Soit $X\in\mathcal{M}$, et soient $\Phi,\Psi$ deux parties de $\mathcal{M}$ (profigures ?) telles que
\[ \Phi \underset{\mathrm{pos}}{<} X < \Psi \]
[NB si $\Phi$ et $\Psi$ sont des profigures, l'existence d'un tel $X$ équivaut à : $\Phi \overset{1}{<} \Psi$ ; si $\operatorname{ex}\delta\Phi = a < \operatorname{or}\delta\Psi = b$, on prend $X=S_{a,b}$. Si $a=b$, on prend $X=P_a$].
Alors
\[ \operatorname{supp}^\circ(\Phi\cup\Psi) = \operatorname{supp}^\circ(\Phi)\cup\operatorname{supp}^\circ(\Psi) \]

\marginal{\uncertain{Ici} il faut \uncertain{supposer} $\Phi,\Psi$ des profigures}

\textbf{Démonstration} On a évidemment $\supset$, il faut prouver $\subset$, i.e.
\[ Z\in\operatorname{supp}^\circ(\Phi\cup\Psi) \overset{?}{\Longrightarrow} Z\in\operatorname{supp}^\circ\Phi \ \text{ou}\ Z\in\operatorname{supp}^\circ\Psi . \]
\struck{\ill{}} Notons que
\[ \Phi \underset{\mathrm{pos}}{<} X \underset{\mathrm{pos}}{<} \Psi \]
\uncertain{implique que}
\[ X\in\operatorname{cosupp}^\circ(\Phi\cup\Psi), \]
donc $Z\,|{\circ}|\,X$, ce qui signifie \ill{}
\[ Z \underset{\mathrm{pos}}{\leq} X \ \text{ou}\ Z \underset{\mathrm{pos}}{\geq} X . \]

\page{113}\note{p. 105 de l'auteur.}
Je vais prouver
\[ Z<X \Longrightarrow Z\in\operatorname{supp}^\circ(\Phi) \]
\[ Z>X \Longrightarrow Z\in\operatorname{supp}^\circ(\Psi) . \]
\struck{Pour ceci,} Prouvons p.ex. la première, (\uncertain{la seconde} \ill{})
\[ Z<X \quad (\text{et } Z\in\operatorname{supp}^\circ(\Phi\cup\Psi)) \]
Soit donc
\[ Z'\in\operatorname{cosupp}^\circ(\Phi), \quad\text{i.e.}\quad Z'\,|{\circ}|\,\Phi \]
nous \struck{\ill{}} devons prouver $Z\,|{\circ}|\,Z'$. Nous \struck{Soit} savons \ill{} \uncertain{que} \add{$Z\,|{\circ}|\,\Psi$, \ill{} \uncertain{qui} si} $Z'\,|{\circ}|\,\Phi\cup\Psi$ (\uncertain{c'est} \ill{} \uncertain{seulement} $Z'\,|{\circ}|\,\Phi$).

\struck{Mais soit} Cas \uncertain{où} « \uncertain{singulier} » $X$, \uncertain{cas} \ill{}
\[ \delta X = \{\operatorname{ex}\delta\Phi,\ \operatorname{or}\delta\Psi\} \quad (\operatorname{ex}\delta\Phi=a,\ \operatorname{or}\delta\Psi=b) \]
i.e.
\[ \begin{cases} X=P_a & \text{si } a=b \\ X=S_{a,b} & \text{si } a<b \end{cases} \]

\struck{\ill{}} Comme $Z'\,|{\circ}|\,\Phi$, il s'ensuit \struck{\ill{}} $Z'\,|{\circ}|\,P_a$ ou $Z'=P_a$, \uncertain{simplement} que $\delta\Phi\cup\delta Z'\in\operatorname{Drap}(\mathcal{L})$, \uncertain{ou} \uncertain{fortiori} (car on aurait $\delta\Phi\cup\{P_a\}\in\operatorname{Drap}\mathcal{L}$, \struck{\ill{}} si on avait $Z'=P_a$, comme par hypothèse $P_a\notin Z'\,|{\circ}|\,\Phi$, et $P_a\,|{\circ}|\,\Psi$ puisque $X<\Psi$, on a $Z'\,|{\circ}|\,\Phi\cup\Psi$, et on a gagné. Sinon on aura $Z' \underset{\mathrm{pos}}{<} P_a$ ou $Z' \underset{\mathrm{pos}}{>} P_a$. Si
\[ Z' \underset{\mathrm{pos}}{<} P_a, \ \text{on aura}\ Z' \underset{\mathrm{pos}}{<} \Psi, \]
donc $Z'\,\|\,\Psi$, on a gagné encore.
\note{la parenthèse ouverte après « \uncertain{fortiori} » n'est pas refermée sur la page.}

\marginal{\textbf{Lemme} (trivial) $Z'\,|{\circ}|\,\Psi$, et $Z'\,|{\circ}|$ \ill{}, alors $Z'=$ \ill{} \ill{} \textbf{Lemme} (trivial) $X \underset{\mathrm{pos}}{<} \Psi$ et $Y \underset{\mathrm{pos}}{\leq} X$ $\Rightarrow$ \ill{}}
\note{deux petits lemmes en marge, écrits en biais et en partie surchargés.}

\page{114}\note{p. 106 de l'auteur.}
Reste le cas \struck{$Z'>P_a$} $Z' > P_a$. \struck{Mais} comme $Z<X$, on a $Z\leq P_a$ ou \struck{$Z<P_a$}, donc on aura $Z<Z'$, d'où a fortiori $Z\,|{\circ}|\,Z'$, cqfd.

NB Je n'ai pas eu à utiliser Hyp., \struck{ni la constructibilité de $\Phi$} mais \struck{\ill{}} j'ai utilisé \struck{le constat} que $\Phi,\Psi$ sont des profigures.

\textbf{Dém. de (b)} Résulte par récurrence \add{sur $\operatorname{card}\delta\Phi$} de 1°, \struck{\ill{}} et du

\textbf{Lemme} : Soient $a<b<c$, et
\[ \Phi = \{S_{a,b},\ P_b,\ S_{b,c}\} . \]
Alors
\[ \operatorname{supp}^\circ(\Phi) = \operatorname{supp}^\circ(S_{a,c}) \]
(Lemme de composition des segments).

\textbf{Démonstration} Il suffit de voir que $\Phi$ et $S_{a,c}$ ont même $\operatorname{cosupp}^\circ$. \struck{Soit} \struck{Cela résulte} \uncertain{de ce} résultat plus précis : Si $X\in\mathcal{M}$, on a
\[ X \underset{\mathrm{pos}}{<} \Phi \ (\text{i.e. } X<S_{a,b}) \iff X<S_{a,c} \]
\[ X \underset{\mathrm{pos}}{>} \Phi \ (\text{i.e. } X>S_{b,c}) \iff X>S_{a,c} \]
trivial (sans Hyp).

\page{115}\note{p. 107 de l'auteur.}
Ainsi, le théorème est établi. On en conclut que \struck{tout}

\note{à partir d'ici, le reste de la page (le corollaire) est biffé de traits obliques ; il est transcrit sans marquer chaque ligne.}
\textbf{Corollaire} (détermination des supports constructibles de $\mathcal{M}$). \struck{Soit $X$}
1) Tout support constructible \uncertain{peut s'obtenir} de la façon suivante : on prend une \struck{\ill{}} figure
\[ F = F_D, \quad D = (T_1 \underset{\mathrm{pos}}{<} T_2 < \cdots \underset{\mathrm{pos}}{<} T_n) \in \operatorname{Drap}^*(\operatorname{Drap}^*(\mathcal{L})), \]
et pour tout $T_i$, on prend une partie
\[ \beta_i \subset \partial F_{T_i} = \partial T_i . \]
\struck{\ill{}} On lui associe la partie constructible
\[ \Phi = (F_D\setminus F_D\cap\mathcal{M}_0)\cup\bigcup\beta_i
= \bigcup_i\big((F_{T_i}\setminus F_{T_i}\cap\mathcal{M}_0)\cup\beta_i\big) \]
et
\[ \operatorname{supp}^\circ\Phi \overset{\text{th}}{=} \operatorname{Omb}^\circ(\Phi) \]
\note{le signe $=$ est peut-être barré ($\neq$).}
\[ \operatorname{supp}^\circ\Phi = \operatorname{supp}^\circ\Phi' \]
$\Phi'$ déduit de $\Phi$ en « recollant » les morceaux », et ayant la propriété d'être « \uncertain{non morcelée} », i.e.
\[ \forall x\in\Phi'\cap\mathcal{M}_0, \quad \operatorname{card}\{X\in\Phi' \mid x\triangleleft X\}\ \text{\struck{$\neq 2$}} \]
(NB ce cardinal est pris $\in\{0,1,2\}$)

Ceci posé, on trouve :

\page{116}\note{p. 108 de l'auteur.}
\textbf{Théorème} 1) Tout support constructible $S$ s'écrit sous la forme $\operatorname{supp}^\circ\Phi$, où $\Phi$ est une profigure non morcelée\add{*} ; \struck{2) Cette écriture est unique, et on a} et pour une telle $\Phi$, on a
\[ \operatorname{supp}^\circ(\Phi) = \operatorname{Omb}^\circ(\Phi) = \bigcup_{X\in\Phi}\operatorname{Omb}^\circ(X) \]
\[ \text{(réunion disjointe).} \]

\marginal{* je dis plutôt ?? « irredondante »}

\struck{$\Phi_1=$} Les éléments de $\Phi\cap\mathcal{M}_1$ sont les éléments 1) de dim $1$ et de $S$ qui sont \struck{maximaux} \add{(ou bien \ill{}, c'est pareil)} 2) maximaux de $S$ pour $\overset{\circ}{<}$. Les \add{\uncertain{constituants}} \struck{éléments de} éléments de $\Phi_0=\Phi\cap\mathcal{M}_0$ sont \struck{les éléments de $S\cap\mathcal{M}_0$ qui sont suivants : où bien ils sont}
\[ \Phi_0 = \Big(\bigcup_{X\in\Phi_1}\partial X\Big)\cap S \cup \Phi_{0,\mathrm{is}} \]
(réunion disjointe), où \struck{$\Phi_{0,\mathrm{is}}$ est}
\[ \Phi_{0,\mathrm{is}} = \text{ens. des sommets isolés de } \Phi \]
\[ = \text{ens. des éléts de } S\cap\mathcal{M}_0 \text{ qui sont maximaux dans } S \text{ pour } \ll \]
\struck{\ill{}}

\struck{\ill{}} En particulier, $\Phi$ est bien déterminé par la connaissance de $S$. On dit que $S$ est fermé, si $\Phi$ est fermée, i.e. que $S$ est le support \uncertain{vrai}, si $S$ est le support d'une figure.

\page{117}\note{p. 109 de l'auteur.}
2) Pour tout support constructible $S$ \struck{\ill{}} \add{$=\operatorname{supp}^\circ\Phi$}, il existe un plus petit support constructible fermé \add{qui le contient}, il est donné par
\[ \overline{S} = \operatorname{supp}^\circ(\overline{\Phi}) = \operatorname{supp}(\overline{\Phi}) \]
— mais je ne suis pas sûr \ill{} que \struck{$=$ si $\Phi$} \ill{} non morcelée, $=$ si $\Phi$ \ill{} $\overline{\Phi}$ peut fort bien être morcelée, si $\Phi$ ne l'est pas.

\marginal{Vérifier ! Facile : c.à.d. aussi $\operatorname{supp}^\circ(\widetilde{S})=\overline{S}$ ? et $=$ adhérence au sens des relations $\triangleleft$ dans $\mathcal{M}$.}
\drawing{en marge : un segment $\Phi$ formé de deux intervalles fermés raccordés en un point $(\bullet)$, « $\Phi$ non morcelée », et en dessous le même avec $\overline{\Phi}$, « morcelée ».}

\struck{3)} 3) Description combinatoire des supports constructibles via les profigures non morcelées ?

\marginal{4. juillet}\note{la date est d'une écriture plus lourde ; le premier chiffre se lit 4, comme en p.~102.}
Je voudrais donner une présentation « autoduale » des supports constructibles \uncertain{tout} par \uncertain{les} profigures non morcelées (ou « irredondantes »). Pour ceci il faut « omettre » \ill{} \uncertain{pour} ses points isolés \add{de $\Phi$}, et les traits ou \uncertain{parties} « points \uncertain{isolés} » i.e. ceux des points de $\mathcal{L}$ qui \struck{\ill{}} ne sont pas dans $\Phi$, mais adhérents : $\Phi$ dans $\mathcal{M}$.

\marginal{\add{et pas $\in\partial\Phi$} — i.e. les « points frontière intérieurs »}

On procède ainsi. \struck{$\Phi_1 =$} \struck{$\overline{\Phi}\setminus$ ens. des points isolés}
On regarde la figure
\[ \text{\struck{$\overline{\Phi}$}}\qquad \Phi_1 = \overline{\Phi}\setminus\{\text{ens. des points isolés de } \overline{\Phi}\} \]
c'est une figure qui \struck{\ill{}} peut ne pas être irredondante. Celle-ci l'est ssi \struck{$\Phi$}
\[ \overline{\Phi}\setminus\Phi \subset \partial\Phi \quad \text{\struck{$\overline{\Phi}\setminus\Phi$}} = \{s\in\delta\Phi \mid \text{ordre de } s \text{ dans } \overline{\Phi} \text{ est } 1\} \]
\struck{est formé des points de}
i.e. ssi $\Phi$ \ill{} pas de point frontière intérieur.

\page{118}\note{p. 110 de l'auteur.}
\textbf{Parenthèse} Zoologie des sommets d'une préfigure.

\drawing{une droite portant, de gauche à droite : un point plein $a$, un trait jusqu'à un point ouvert $(\ )$ en $b$, un trait jusqu'à un point plein entre parenthèses $(\bullet)$ en $c$, un trait jusqu'à un crochet en $d$, puis un point isolé $e$.}
\[ \Phi = \{P_a,\ S_{a,b},\ S_{b,c},\ P_c,\ S_{c,d},\ e\} \]
\[ \overline{\Phi} = \{P_a,\ S_{a,b},\ P_b,\ S_{b,c},\ P_c,\ S_{c,d},\ d,\ e\} \]
\[ \partial\Phi = \{a,\ d\} \]
\note{plusieurs termes de ces deux ensembles sont surchargés ; la lecture des $P$ est incertaine.}

\note{à droite, séparée par un trait vertical, une première rédaction biffée de traits obliques, reprise plus bas : « $\delta\Phi$ (ens. des sommets de $\Phi$) $=\overline{\Phi}\cap\mathcal{L}=(\overline{\Phi})_0$ ; sommets propres : $\delta\Phi\cap\Phi_0$, impropres $\delta\Phi\cap\mathcal{L}$ ; $\operatorname{ord}(a,\Phi)=\operatorname{card}(\{X\in\mathcal{M}_1\cap\Phi=\Phi_1 \mid a\triangleleft X\})$ ; $\delta\Phi=\{a\in\mathcal{L} \mid \operatorname{ord}(a,\Phi)=\}$ ; sommet isolé de $\Phi$ : d'ordre $0$ ; sommet bord de $\Phi$ : d'ordre $1$, deux cas : $a\in\Phi$ sommet bord propre, $a\notin\Phi$ sommet bord impropre ».}

Reprenons
\[ \delta\Phi \ (\text{ensemble des sommets de } \Phi) = \overline{\Phi}\cap\mathcal{L} = (\overline{\Phi})_0 \]
sommets propres : ceux de $\Phi_0=\overline{\Phi}\cap\mathcal{L}\subset\delta\Phi$ ;
sommets impropres : ceux qui ne sont pas dans $\Phi_0$
\[ \operatorname{ord}(a,\Phi) = \operatorname{card}(\{X\in\Phi_1 \mid a\triangleleft X\}) \in \{0,1,2\} \]
\note{sous $a$ : « $\in\mathcal{L}$ » ; sous $\Phi_1$ : « $=\mathcal{M}_1\cap\Phi$ ».}
(ordre de $a$ p.r. à $\Phi$) Cet ordre est zéro si $a$ n'est pas un sommet, mais il peut être zéro aussi pour un sommet de $\Phi$, \struck{\ill{}}

\textbf{Lemme} Si $a\in\delta\Phi$, \struck{alors} \add{et} $\operatorname{ord}(a,\Phi)=0$, on a $a\in\Phi$. Un tel sommet de $\Phi$ est dit \struck{isolément} isolé.
\[ \Phi_{\mathrm{is}} = \{s\in\delta\Phi \mid \operatorname{ord}(s,\Phi)=0\} \subset \Phi_0 \]
Les sommets d'ordre $1$ sont les \emph{sommets-bord}
\[ \partial\Phi = \{s\in\delta\Phi \ (\text{ou } s\in\mathcal{L}) \mid \operatorname{ord}(s,\Phi)=1\} \]
Parmi les sommets-bord, on distingue les sommets-bord propres, et impropres
\[ \Phi_0\cap\partial\Phi = \partial_0\Phi \subset \Phi \]
\[ \text{(sommets-bord propres)} \]
Les sommets d'ordre $2$ sont les \emph{sommets internes}
\[ \delta_{\mathrm{int}}(\Phi) = \{s\in\delta\Phi \ (\text{ou } s\in\mathcal{L}) \mid \operatorname{ord}(s,\Phi)=2\} \]

\page{119}\note{p. 111 de l'auteur.}
Parmi les sommets internes, on distingue les sommets propres, ou sommets \emph{redondants} ou « sommets de morcellement », et les sommets impropres, ou \emph{sommets lacunaires}.

Dans la profigure
\drawing{la même figure qu'en p.~118 : un point plein $a$, un point ouvert en $b$, un point plein entre parenthèses en $c$, un crochet en $d$, un point isolé $e$.}
il y a un sommet de chacun des cinq types combinatoires
\begin{itemize}
\item sommet isolé $e$
\item sommet bord propre $a$, impropre $d$
\item sommet interne propre ou redondant $c$
\item sommet interne impropre ou lacunaire $b$
\end{itemize}

\struck{Pour une intervalle} J'\uncertain{aime} \uncertain{mieux} : \uncertain{préférer} appeler $S_{a,b}=S_{b,a}$ ou $S_{\{a,b\}}$ un \emph{intervalle} (plutôt qu'un segment). \struck{Le} \uncertain{Le} \emph{segment} associé à l'intervalle est
\[ \widetilde{S}_{a,b} = \{S_{a,b},\ a,\ b\} \]
— c'est le \uncertain{multisimplexe} complet associé : \ill{}.

\struck{Si $X\in\Phi_{T,D}$} On a posé
\[ \Phi_0 = \Phi\cap\mathcal{M}_0 = \Phi\cap\mathcal{L}, \qquad \Phi_1 = \Phi\cap\mathcal{M}_1 \]
\note{sous $\Phi_0$ : « ou ens. des sommets propres ».}

\note{le paragraphe qui suit est biffé de traits obliques.}
Pour les éléments de $\Phi_1$, ou \emph{intervalles de $\Phi$}, on distingue les intervalles de $\Phi$ tels que $\widetilde{X}\subset\Phi$, i.e. $\overline{X}\subset\Phi$, i.e. $\partial X\subset\Phi_0$ (\uncertain{ou} \emph{fermés dans $\Phi$}), et ceux pour lesquels $\partial X\cap\Phi_0=\{a\}$ (\uncertain{semi}-\emph{ouverts} dans

\page{120}\note{p. 112 de l'auteur.}
Regardons l'opération suivante sur $\Phi$
\[ \Phi \longmapsto \overline{\Phi} \longmapsto \overline{\Phi}\setminus\overline{\Phi}_{\mathrm{is}} \Longrightarrow \Phi_{\mathrm{comp}} = \Psi \]
\note{sous $\overline{\Phi}$ : « figure \struck{\ill{}} associée à $\Phi$, ou adhérence » ; sous $\overline{\Phi}_{\mathrm{is}}$ : « $=\Phi_{\mathrm{is}}$ » ; sous $\overline{\Phi}\setminus\overline{\Phi}_{\mathrm{is}}$ : « adhérence purement $1$-dimensionnelle » ; sous $\Phi_{\mathrm{comp}}$ : « ou “\struck{télescopage}\add{complétion} par effacement des pts redondants de $\overline{\Phi}\setminus\overline{\Phi}_{\mathrm{is}}$” ». Après $\Phi_{\mathrm{comp}}$, quelques mots biffés.}

On trouve une \emph{figure} dont les composantes connexes sont des segments, donc donnée par une \struck{suite} drapeau dans $\mathcal{M}_1$, ordonné par $\overset{1}{<}$
\[ D \in \operatorname{Drap}^*(\mathcal{M}_1,\overset{1}{<}) \]
\[ D = \{X_1,\ X_2,\ \ldots,\ X_n\} \]
\[ X_i = S_{a_i,b_i}, \]
\drawing{une droite portant les segments $[a_1,b_1]$, $[a_2,b_2]$, $[a_3,b_3]$, $[a_n,b_n]$, marqués de croix à l'intérieur et entre eux ; le dernier segment est tracé en gras.}
de telle façon que
\[ \begin{cases}
\Psi = \Phi_{\mathrm{comp}} = \Phi_D = \Phi_{X_1,X_2,\ldots,X_n} \\
\delta\Psi = \partial\Psi = \partial\Phi = \{a_1,b_1,a_2,b_2,\ldots,a_n,b_n\} .
\end{cases} \]

Pour \struck{en donner} \add{décrire} \uncertain{l'écriture} $\Phi$ en termes de $\Phi_{\mathrm{comp}}$ il faut :

1°) Donner
\[ \delta\Phi \in \operatorname{Drap}^*(\mathcal{L}) \supset \delta\Psi \]
drapeau de $\mathcal{L}$ contenant $\delta\Psi$. Les nouveaux pts rajoutés\add{*} sont \uncertain{donc} \struck{points} les \add{sommets} isolés et les sommets internes de $\Phi$.

\marginal{* \uncertain{qui} \uncertain{sont} \ill{} \ill{}}

\note{la page s'arrête sur le point 1° ; la suite de l'énumération continue au-delà de ce lot.}

\end{document}
